\chapter{How to read these drawings}
\label{ch:method}

\section{A book as a walk}

Imagine reading a novel aloud while walking across an empty field. You take one
step for every word. At the end of each sentence you stop and ask how it felt.
If it was warm, hopeful or kind, you turn a quarter turn to the left. If it was
cold, frightened or cruel, you turn a quarter turn to the right. If it was
neither, a plain sentence about a door or a staircase, you keep walking
straight ahead. Then you read the next sentence.

The drawings in this book are the footprints of that walk. Every sentence is a
straight line as long as the sentence has words, and every corner is a change
of mood. Nothing has been placed by hand. The shape of a drawing is decided
entirely by the order of the sentences, their lengths, and how they feel.

How a sentence ``feels'' is measured with a published word list, AFINN, which
gives some 3,400 English words a score from $-5$, reserved for the worst
insults, through $-4$ for \emph{torture}, to $+5$ for \emph{superb} or
\emph{breathtaking}. A sentence's score is the sum of its words' scores. Most words
are not on the list at all, which is why so many sentences walk straight on.

Names are not feelings, so a capitalised word in the middle of a sentence, like
\emph{Miss} in Miss Bürstner or \emph{Block}, the businessman in \emph{The
Trial}, is not scored, and neither is any word a book mainly uses as a name.

Sentences are numbered in reading order, and the book refers to them that way:
\sref{1} is the first sentence of the novel.

\section{Seeing past chance}

Any long walk of this kind will tangle. Lines will cross, loops will form,
some regions will turn black with ink. The question is always which of these
shapes belong to \emph{this} book, and which any book would make.

So every drawing is compared with its own shadows. The book's sentences are
shuffled like a deck of cards, two hundred times, and each shuffled version is
walked and drawn. A shuffled book has exactly the same sentences, with the same
lengths and the same feelings. Only the order is lost. When the real book
behaves like its shuffles, the shape is simply what its material makes. When
it behaves as none of its shuffles do, the shape is the author's, because it
comes from the order in which things are told.

The essays in this book only call something a property of a novel when it
stands clearly apart from its shuffles. The figures behind each such claim are
collected in the notes at the end of the book.

\section{Shapes to look for}

\begin{description}
\item[Knots] places where the path crosses itself again and again. A knot is
most interesting when the sentences that make it come from far-apart parts of
the book: scenes that the author keeps apart, which the drawing lays on top of
each other.
\item[The hub] the single sentence that is crossed most often by others.
\item[The return] the crossing between the two sentences that lie farthest
apart in the book, where the late book walks across its own beginning.
\item[Coils] sentences that keep turning the same way spiral tightly. A coil is
a passage that dwells in one feeling.
\item[Corridors] runs of plain sentences draw long straight lines: description,
procedure, a room being looked at.
\item[The darkest passage] the thirty sentences in a row with the lowest score.
\item[Palimpsests] sometimes a later sentence is drawn exactly on top of an
earlier one, the way old manuscripts were written over.
\end{description}

\section{The mineshaft}

The same walk can also be built in three dimensions. As before, each warm
sentence turns the path left and each cold one turns it right, but now every
left turn also tilts the path slightly upwards and every right turn tilts it
slightly down. A book that turns right more than left sinks; a book that
lingers in one feeling tunnels steeply. Seen from the side, the book becomes a
shaft dug through the ground, with the chapters marked along it.

\section{More instruments}

The walk is one way of looking. The later chapters use five more, each a
different picture of the same text:

\begin{description}
\item[The vocabulary record] the whole dictionary laid out on a disc like the
groove of a record, from the commonest word at the centre to the rarest at the
rim. Every word a book uses is a black dot.
\item[Tree rings] the book cut like a trunk: the first page at the centre,
each chapter a ring whose area is its length, each sentence a grain as long as
its words are long.
\item[The loom] the names a book uses most are the warp threads, and the book
runs down through them as weft. A black stitch marks each stretch where a name
occurs, so the whole cast becomes cloth.
\item[Echoes] every recurring word is drawn as one arch over the length of
the book, spanning its longest absence. High arches are words that return from
far away.
\item[The tide] the share of words spoken aloud, inside quotation marks, as a
black tide rising and falling through the book.
\end{description}

\section{What the drawings cannot see}

A word list is a blunt instrument. It does not hear irony, it does not notice
``not'' (``not happy'' counts as happy), and it cannot tell a feeling from a
figure of speech. A lawyer's empty reassurances score as hope. For a translated
novel, it reads the translator's words, not the author's.

These blind spots are the same for a book and for its shuffles, so they cannot
by themselves make a book stand apart from chance. But they decide what the
drawings can honestly be said to show: the pattern of the surface of the words,
in the order the author chose. Where a single ambiguous word changes a result,
the essays say so.
